\subsection[Challenge 56: Dark matter detection with Cosmic Explorer]{Challenge 56: Dark matter detection with Cosmic Explorer}
\textbf{Problem ID:} \texttt{Challenge\_56\_main}\\
\textbf{Primary inferred discipline:} Gravitation, Cosmology \& Astrophysics\\
\textbf{Secondary inferred disciplines:} High Energy \& Nuclear Physics\\
\textbf{Python implementation:} \texttt{challenges/56/answer.py}

\subsubsection*{Problem}
\paragraph*{Problem setup}
Consider a scalar field $\phi$ linearly coupled to the standard model electromagnetic field with coupling strength $\Lambda_\gamma^{-1}$. We can write the interaction term of the Lagrangian as $\mathcal{L}_\text{int}=\frac{\phi}{4 \Lambda_\gamma}F^{\mu\nu}F_{\mu\nu}$. Assuming the scalar field makes up all the dark matter, it may induce observable effects in laser interferometers. We consider signals of scalar field in the Cosmic Explorer, whose planned strain sensitivity is about $2\times 10^{-25}\, \text{Hz}^{-1/2}$ in the frequency range of $50\,\text{Hz} \leq f \leq 500\,\text{Hz}$. The interferometer arm length of Cosmic Explorer is designed to be $40\, \text{km}$. We assume the Cosmic Explorer uses a similar beamsplitter as LIGO, with thickness $6\,\text{cm}$ and index of refraction $3.5$.

\paragraph*{Main problem}

If there is an overdensity of dark matter at the device that is $178$ times the local dark matter density (take the local dark matter density to be $0.4\,\text{GeV}/\text{cm}^3$ and velocity to be $230\,\text{km}/\text{s}$), what is the smallest $\Lambda_\gamma^{-1}$ the Cosmic Explorer can probe at $200\,\text{Hz}$ with a signal-to-noise ratio of 1 and with an observation time of $1000\,\text{s}$ or $0.7\,\text{yrs}$, respectively? Keep at least three significant digits in the final result.


\paragraph*{Reference-model assumptions}
Use the simple-Michelson response with no arm-cavity or signal-recycling gain.
Model a freely suspended, centered, 45-degree silicon beamsplitter at 295 K
with a 2-micrometre laser, retaining the specified $n=3.5$ and $l=6$ cm.
Include Snell refraction, the motion of both surfaces about the plate center,
and finite arm light-travel times. Fix the photon-only response by
$\delta m_e=0$, $\delta l/l=-\delta\alpha/\alpha$ and
$\delta n/n=-0.0442\,\delta\alpha/\alpha$; omit relativistic lattice corrections.
Use both halo dispersion and streaming speed equal to 230 km/s, held constant,
and keep the stated overdensity fixed. Treat the given sensitivity as a
one-sided strain amplitude spectral density. With $\omega=2\pi(200\,\mathrm{Hz})$,
$v=230$ km/s, $\tau_c=[(v/c)^2\omega]^{-1}$, define the correlation envelope
and the amplitude threshold at SNR one by
\[
A(\tau)=\frac{\exp[-x^2/(2(1+x^2))]}{(1+x^2)^{3/4}},\quad x=\tau/\tau_c,
\qquad h_{\min}(T)=\sqrt{S_h/T}\left[
\frac{T/3}{\int_0^T(1-\tau/T)^2A(\tau)^2\,d\tau}\right]^{1/4}.
\]
Use Julian years and report $\Lambda_\gamma^{-1}$ in $\mathrm{GeV}^{-1}$.
The target is the primary finite-coherence reference calculation, not the
separate sharp two-regime approximation or a forecast for an unspecified
full Cosmic Explorer optical configuration.
